
\subsection{Open-chain gap theorem and boundary remainder}
\label{obc:setup}

We extend the diagrammatic doubling argument from periodic rings
to open chains with endpoint magnetic fields. Consider the spin-1
Heisenberg chain with equal endpoint fields (\fieldsame),
\begin{equation}
    H_n^{\same}
    =\sum_{i=1}^{n-1}\bm S_i\cdot\bm S_{i+1}
    -h(S_1^z+S_n^z),\qquad h=\frac35,
    \label{obc:model}
\end{equation}
where $\sum_\alpha(S^\alpha)^2=2I$ and the exchange coupling is one.
We give a diagrammatic proof of the following theorem from
Ref.~\cite{OpenAI2026Boundaryfield}, using its certified finite inputs.

\begin{theorem}[Boundary-field gap~\cite{OpenAI2026Boundaryfield}%
        ~\leanmark{https://sirui-lu.com/CertifiedGaps/docs/CertifiedGaps/Heisenberg/OpenChain.html\#CertifiedGaps.SpinOneHeisenberg.spin_one_boundary_field_gap}]
    \label{obc:theorem}
    For every odd $n\ge1921$, $H_n^{\same}$ has a unique ground state and
    \[
        \Delta_n>\frac{\log 10}{392}\simeq0.005874.
    \]
\end{theorem}
\noindent The two constants come from three doublings of the certified
seed $(M_2,\beta_2)=(240,49)$; scale $5$ is the first at which every
tracked bound is below $10^{-2}$:
\[
    M_5=2^3\cdot240=1920,\qquad \beta_5=2^3\cdot49=392,\qquad
    n\ge M_5+1=1921.
\]

Without the endpoint fields, the effective spin-$1/2$ modes at the
two ends produce four low-lying states whose splitting is exponentially
small in the chain length~\cite{Kennedy1990Exact}. Applying equal fields on an odd chain lifts this
boundary degeneracy~\cite{OpenAI2026Boundaryfield}.

Our strategy is to bound the thermal purity of the open chain at a
scale $\beta$ proportional to the length, just as in the periodic
argument. Order the Gibbs weights of $H_n^{\same}$ as $w_0\ge w_1\ge\cdots$.
The thermal purity $\Pt^{\same}(n,\beta)=\sum_m w_m^2$ then bounds the
leading weight and the gap:
\begin{equation}
    w_0\ge\Pt^{\same}(n,\beta),\qquad
    e^{-\beta\Delta_n}
    \le\frac{1-\Pt^{\same}(n,\beta)}{\Pt^{\same}(n,\beta)},\qquad
    \Delta_n=E_1(n)-E_0(n).
    \label{obc:target}
\end{equation}
A thermal purity above $1/2$ forces a unique ground state and bounds the
gap from below. What changes in the open-chain setting is how we
propagate the thermal purity across doubled lengths and inverse
temperatures.

To see where the new difficulty enters, consider the spatial transfer
operator at a fixed temperature. Suppose that $X_\beta$ has a simple
positive leading eigenvalue $\lambda$ (\specbars[mpored]{.9,.06,.04}), and
write the normalized transfer operator as
\begin{equation}
    T:=X_\beta/\lambda=\Pi+R,\qquad
    \Pi=\ket{\psi}\bra{\psi},\qquad
    \Pi R=R\Pi=0,\quad \|\psi\|=1.
    \label{obc:split}
\end{equation}
In a periodic ring, the trace weights the leading projector $\Pi$ by one.
In an open chain, the chain ends are contracted with boundary vectors,
so the leading contribution is weighted instead by how much of $v$ lies
along $\psi$ (\overlapglyph):
\begin{equation}
    \underbrace{\tnchain[ring]{\Pi}}_{\trv\Pi=1},
    \qquad
    \underbrace{\tnchain[capl=\bar v,capr=v]{\Pi}}_
    {c=|\langle\psi,v\rangle|^2}.
    \label{obc:leading-diagram}
\end{equation}
Because $c$ is not fixed by normalization and could a priori vanish,
squaring in space no longer closes on its own. With $T$ self-adjoint, the
same boundary vector at both ends, and an even length $m$, the remainder of the
normalized boundary matrix element is a squared norm, hence nonnegative:
\begin{align}
    B_m
        &=\tnchain[capl=\bar v,capr=v]{T^m}\notag\\
        &=\tnchain[capl=\bar v,capr=v]{\Pi}
    +\tnchain[capl=\bar v,capr=v]{R^m}
    =c+D_m,\notag\\
    D_m &=\|R^{m/2}v\|^2\ge0.
    \label{obc:boundary-split}
\end{align}
To show that the leading eigenvector dominates $B_m$, we bound the
remainder's share of $B_m$ rather than the remainder itself,
\begin{equation}
    e_m:=\frac{D_m}{B_m}=1-\frac{c}{B_m},
    \label{obc:defect}
\end{equation}
the grey fraction of \stackbar{.14} at the certified seed value
$e_{142}\le.14$ of Eq.~\eqref{obc:seeds}. Increasing the chain length
suppresses $D_m$ by powers of $\|R\|$ while leaving $c$ unchanged.

Doubling the inverse temperature changes both $v$ and $\psi$, so $c$
does not bound the new overlap $c'$; Step~1 (Sec.~\ref{obc:move-one})
removes $c'$ by comparing two lengths at the new temperature.

The recurrence therefore tracks three contractions at every scale $j$,
each with its small quantity and its certified seed value:
\begin{center}
    \renewcommand{\arraystretch}{1.5}
    \begin{tabular}{@{}clllc@{}}
        \toprule
            & contraction                        & tracked quantity                & seed ($j=2$) & controlled in                                  \\
        \midrule
        (R) & \obcring\ ring purity              & $\varepsilon_j\ge1-\Ps,\,1-\Pt$ & $.052$       & Sec.~\ref{obc:ring}                            \\
        (E) & \obcends\ boundary overlap         & $e_{k_j}(j)=D/B$                & $.14$        & Sec.~\ref{obc:move-one}                        \\
        (T) & \obcthermal\ opposite-field purity & $1-\Pt^{\opp}(M_j,\beta_j)$     & $.21$        & Secs.~\ref{obc:move-two}--\ref{obc:move-three} \\
        \bottomrule
    \end{tabular}
\end{center}
Section~\ref{obc:caps} constructs the boundary vectors, Sec.~\ref{obc:initialization}
runs three explicit steps from $j=2$ to $j=5$ and closes the induction,
and Sec.~\ref{obc:comparison} passes to equal fields and odd lengths.

\subsection{Continuous-time transfer kernel and boundary vectors}
\label{obc:caps}

We first represent the open-chain partition functions as transfer-operator
matrix elements using the continuous-time bond expansion of
Ref.~\cite{OpenAI2026Boundaryfield}, which sums over the times of bond
insertions rather than placing them on a finite Trotter grid.
As in the periodic proof (Sec.~\ref{cert:sec:transfer}), we shift the
Hamiltonian by a constant $\kappa$ per site. The shift is chosen so that $-\kappa$
lies just above the energy per site of the $72$-site matrix-product
trial state of Ref.~\cite{OpenAI2026Periodic} ($E_0(72)/72<-\kappa$,
Sec.~\ref{cert:sec:inputs}):
\begin{equation}
    \kappa=\frac{700741}{500000}=1.401482.
    \label{obc:kappa}
\end{equation}
Throughout this section, $Z_n$ and all open-chain partition functions
include the shift $\kappa nI$. It places the leading transfer eigenvalue
near one at the seed scales and cancels from every purity and spectral
gap.

In the Cartesian basis, $G_\alpha=iS^\alpha$ is real antisymmetric and
\[
    -\bm S_i\cdot\bm S_{i+1}
    =\sum_{\alpha=x,y,z}G_{\alpha,i}G_{\alpha,i+1}.
\]
Expanding $e^{-\beta H}$ on each bond produces a bond history $a$, an
ordered list of events $(\tau,\alpha)$ with $0<\tau<\beta$ and
$\alpha\in\{x,y,z\}$. The bond measure $\mu_\beta$ sums the labels and
integrates the ordered times, with unit mass on the empty history:
\[
    \mu_\beta(\text{all histories})
    =\sum_{k\ge0}\underbrace{3^k}_{\text{labels}}\,
    \underbrace{\frac{\beta^k}{k!}}_{0<\tau_1<\cdots<\tau_k<\beta}
    =e^{3\beta}.
\]
Between the left bond history $a$ and the right bond history $b$ of a
site, the shifted local tensor is
\begin{equation}
    W_\beta^{ss'}(a,b)
    =e^{-\kappa\beta}
    \bra{s}\mathcal T_\tau
    \prod_{(\tau,\alpha)\in a\cup b}G_\alpha\ket{s'},
    \label{obc:local}
\end{equation}
where $\mathcal T_\tau$ orders the events of both bonds by time, later
ones to the left:
\begin{proofdiagram}
    \worldlinekernel
\end{proofdiagram}
Tracing the physical indices defines the transfer kernel $X_\beta(a,b)$
on $L^2(\mu_\beta)$. It is real, and symmetric because the product
follows the times, not the sides.

To represent a boundary field along $z$, we contract the endpoint bond
history with the boundary vector
\begin{equation}
    v_\beta(a)=
    \begin{cases}
        (ih)^{|a|}, & \text{all event labels are }z, \\
        0,          & \text{otherwise},
    \end{cases}
    \qquad \|v_\beta\|^2=e^{\beta h^2}.
    \label{obc:vector}
\end{equation}
Its value on the empty history is one, and its norm follows by summing
$(\beta h^2)^k/k!$ over the number $k$ of $z$ events.

Each event on the boundary bond inserts $G_z=iS^z$ at the end site at
its time. It enters the time-ordered product like an interior bond
event, and summing counts and integrating times exponentiates the
endpoint field. Conjugation reverses the sign of $ih$, so contracting
with $v_\beta$ at both ends gives the auxiliary opposite-field chain
\fieldopp\ (left). The $\pi$ rotation $P=\operatorname{diag}(1,-1,-1)$
about $x$ acts on the right boundary vector as $U_P$; it negates every
$z$ event there and restores equal fields \fieldsame\ (right):
\begin{equation}
    \capfields{0}\qquad\qquad\capfields{1}
    \label{obc:capfields}
\end{equation}
\begin{equation}
    H_n^{\opp}=H_n^{\same}+2hS_n^z.
    \label{obc:opposite}
\end{equation}
The ring bounds of Sec.~\ref{obc:ring} concern the spectrum of $X_\beta$,
so both open-chain partition functions must become matrix elements of
it. In diagram form,
\begin{align}
    e^{-\beta(H_n^{\opp}+\kappa nI)} &=\obcW,\notag\\
    Y_n(\beta)
    :=\Tr e^{-\beta(H_n^{\opp}+\kappa nI)}
                                     &=\obccap{X_\beta^n}{v_\beta},\notag\\
    Y_n^{\same}(\beta)
    :=\Tr e^{-\beta(H_n^{\same}+\kappa nI)}
                                     &=\obccap{X_\beta^n}{U_Pv_\beta}.
    \label{obc:capping}
\end{align}
The two end bond histories meet the boundary vectors, and the physical
trace at every site turns the row of tensors into $X_\beta^n$; the two
chains differ only by $U_P$ at the right end.

Spin rotations act on bond histories by real orthogonal operators that
commute with $X_\beta$, so $[U_P,X_\beta]=0$. Contracting the
time-ordered expansion site by site is exact by Fubini's theorem, since
$\|G_\alpha\|=1$ bounds the absolutely summed series by
$3^ne^{3\beta(n-1)+2\beta h}$ before the energy shift.

The recurrence runs on the auxiliary opposite-field chain at even
length $m$, not on the odd equal-field chain of the theorem. With the
same boundary vector at both ends, the remainder of $Y_m(\beta)$ is the
squared norm $D_m\ge0$ of Eq.~\eqref{obc:boundary-split}, on which the
doubling closes.

\subsection{Bulk ring purities and transfer-operator bounds}
\label{obc:ring}

The open-chain recurrence of Steps~1--3 needs three facts about the bulk
transfer operator at every scale $j$. They are a simple positive leading
eigenvalue that $U_P$ fixes, a decay rate for the remainder across half
a chain, and the ratio $r_j$ of leading eigenvalues at consecutive
temperatures. None is certified directly at large $j$, but the periodic
argument proves ring purity bounds at every scale from one certified
seed, and this subsection extracts the three facts from those purities.
The open chain thus borrows its bulk control from the ring; only the
boundary needs new inputs.

The recurrence runs over the dyadic scales
\begin{equation}
    \beta_j=49\,2^{j-2},\qquad n_j=120\,2^{j-2},\qquad
    M_j=2n_j,\qquad j\ge2,
    \label{obc:scales}
\end{equation}
starting at the certified point $j=2$ of Ref.~\cite{OpenAI2026Boundaryfield}.
The ring length controls the bulk and the open length carries the
thermal purity. The boundary ratio $e_k$ is tracked at the certified length
$142$ first and at the previous open length thereafter,
\begin{equation}
    k_2=142,\qquad k_{j+1}=M_j,
    \label{obc:k}
\end{equation}
so $k_j=M_j/2$ for $j\ge3$:
\begin{proofdiagram}
    \obcscales
\end{proofdiagram}
Ref.~\cite{OpenAI2026Boundaryfield} certifies four inputs at $j=2$%
~\leanmark{https://sirui-lu.com/CertifiedGaps/docs/CertifiedGaps/Heisenberg/BoundaryInitTable.html\#CertifiedGaps.SpinOneHeisenberg.BoundaryInit.boundaryDefect_le_of_table}:
\begin{equation}
    \renewcommand{\arraystretch}{1.4}
    \begin{tabular}{@{}llcl@{}}
        \toprule
        input                      & controls         & weights                        & bound                       \\
        \midrule
        $\Ps(120,49)\ge.948$       & bulk spectrum    & \specbars[mpored]{.95,.03,.02} & $\varepsilon_2=1-.948=.052$ \\
        $\Pt(240,49)\ge.948$       & bulk spectrum    & \specbars{.95,.03,.02}         & $\varepsilon_2$             \\
        $e_{142}(2)\le.14$         & boundary overlap & \stackbar{.14}                 & $E_2=.14$                   \\
        $\Pt^{\opp}(240,49)\ge.79$ & thermal          & \specbars{.79,.13,.08}         & $\eta_2=1-.79=.21$          \\
        \bottomrule
    \end{tabular}
    \label{obc:seeds}
\end{equation}
Here $\Ps(n,\beta)=Z_{2n}(\beta)/Z_n(\beta)^2$ and
$\Pt(n,\beta)=Z_n(2\beta)/Z_n(\beta)^2$ are ring purities, $e_m(j)$ is
Eq.~\eqref{obc:defect} at $\beta_j$, and
$\Pt^{\opp}(m,\beta)=Y_m(2\beta)/Y_m(\beta)^2$.

The periodic argument squares $1-\Ps$ and $1-\Pt$ with constant $5$ only
below the threshold $1/100$, which the seed $\varepsilon_2=.052$
exceeds; repeating the four-step update with constant $6$ gives
\begin{equation}
    \Ps(n_j,\beta_j),\ \Pt(M_j,\beta_j)\ge1-\varepsilon_j,\qquad
    \varepsilon_2=.052,\qquad
    \varepsilon_{j+1}=6\varepsilon_j^2.
    \label{obc:ring-purities}
\end{equation}
One scale of the periodic argument, with $\varepsilon=\varepsilon_j$,
updates the spatial and thermal purities in four moves:
\begin{proofdiagram}
    \ringsteps
\end{proofdiagram}
Squaring and renormalizing a positive
trace-one operator of purity $P$ concentrates its weights
($\specbars[black!55]{.5,.3,.2}\mapsto\specbars[black!55]{.66,.24,.1}$),
and its new purity $P'$ satisfies
\begin{equation}
    1-P'\le f(1-P),\qquad
    f(z)=\frac{z^2}{2(1-z)^2}.
    \label{obc:square}
\end{equation}
For $0\le\varepsilon\le.052$, Move~2 gives
\begin{equation}
    f(1-(1-\varepsilon)^3)
    =\frac{\varepsilon^2}{2}
    \left[\sum_{s=1}^3(1-\varepsilon)^{-s}\right]^2
    \le6\varepsilon^2\le\varepsilon,
    \label{obc:six}
\end{equation}
the bracket being at most $3.35$ at $\varepsilon=.052$, so the prefactor is
at most $5.6<6$. Since
$(1-\varepsilon)^2(1-6\varepsilon^2)\ge(1-\varepsilon)^3$, Move~4
gives the same bound $6\varepsilon^2$, and induction on $j\ge2$ proves
Eq.~\eqref{obc:ring-purities}.

Write $\lambda_j$ for the leading eigenvalue of $X_{\beta_j}$ and
$T_j=X_{\beta_j}/\lambda_j=\Pi_j+R_j$ as in Eq.~\eqref{obc:split}. The
three bulk facts, each with the step that consumes it, are%
~\leanmark{https://sirui-lu.com/CertifiedGaps/docs/CertifiedGaps/Open/Criterion.html\#CertifiedGaps.OpenMomentWitness.mu_pow_ge}
\begin{align}
     &\text{(i)}\quad \lambda_j>0\text{ simple},\quad U_P\Pi_j=\Pi_j
     &                                                                           &\text{overlap $c$ in Step~1 and Sec.~\ref{obc:comparison}},\notag\\
     &\text{(ii)}\quad \|R_j\|^{n_j}\le d(\varepsilon_j)<1,\quad
    d(\varepsilon)=\frac{\varepsilon}{2-3\varepsilon}
     &                                                                           &\text{suppression factor $q$ in Step~1},\notag\\
     &\text{(iii)}\quad r_j^{M_j}\ge(1-\varepsilon_j)(1-\varepsilon_{j+1}),\quad
    r_j=\frac{\lambda_{j+1}}{\lambda_j^2}\le1
     &                                                                           &\text{bulk factor in Step~2 and Sec.~\ref{obc:comparison}}.
    \label{obc:ring-bounds}
\end{align}

No transfer eigenvalue is computed, so (i) and (ii) must follow from
the ring purity alone. A purity at least $1-\varepsilon>1/2$ forces one
weight $p_0$ above $1/2$ (\dominantglyph), and its tail bounds every
subleading ratio. The purity is at most $p_0$,
so the tail $z=1-p_0$ is below $1/2$. The purity is also at most
$p_0^2+z^2=1-2z(1-z)$, so $2z(1-z)\le\varepsilon$ and
\[
    z\le\varepsilon,\qquad
    z\le\frac{\varepsilon}{2(1-\varepsilon)},\qquad
    \frac{z}{1-z}\le\frac{\varepsilon}{2(1-\varepsilon)-\varepsilon}=d(\varepsilon).
\]
Applying this estimate to the normalized weights
$p_i=|\lambda_i(\beta_j)|^{n_j}/Z_{n_j}(\beta_j)$ yields the subleading tail bound
\begin{equation}
    \sum_{i\ge1}
    \left|\frac{\lambda_i(\beta_j)}{\lambda_j}\right|^{n_j}
    \le d(\varepsilon_j).
    \label{obc:full-tail}
\end{equation}
Since $d(\varepsilon_j)<1$, three consequences follow.
\begin{itemize}
    \item $\lambda_j$ is simple, and $\|R_j\|^{n_j}\le d(\varepsilon_j)<1$
          is property (ii).
    \item $\lambda_j>0$. A negative leading eigenvalue would make
          $\trv X_{\beta_j}^N<0$ for all large odd $N$, because $X_{\beta_j}$ is
          Hilbert--Schmidt and the rest of its spectrum has smaller radius;
          but $Z_N(\beta_j)>0$.
    \item $U_P\Pi_j=\Pi_j$. The cyclic axis permutation and the three $\pi$
          rotations commute with the real kernel $X_{\beta_j}$, so each acts by a
          sign on the leading eigenline. The permutation has order three, so its
          sign is $+1$. It conjugates the three rotations into one another, so
          their signs agree; two of them multiply to the third, so the common
          sign is $+1$.
\end{itemize}

Property (iii) bounds the bulk factor that Step~2 loses when the
temperature doubles. The two transfer operators do not commute, so we
compare both leading eigenvalues with one ring partition function.
First, $\Pt(N,\beta_j)\le1$ and even $N$th roots as $N\to\infty$
give $\lambda_{j+1}\le\lambda_j^2$, so $r_j\le1$. At length
$M_j=n_{j+1}$,
\begin{align}
    \lambda_{j+1}^{M_j}
     &\ge(1-\varepsilon_{j+1})Z_{M_j}(\beta_{j+1})
     &                                             &\text{leading weight at }\beta_{j+1},\notag\\
     &\ge(1-\varepsilon_{j+1})(1-\varepsilon_j)
    Z_{M_j}(\beta_j)^2
     &                                             &\text{ring thermal purity at }\beta_j,\notag\\
     &\ge(1-\varepsilon_{j+1})(1-\varepsilon_j)
    \lambda_j^{2M_j}
     &                                             &\text{even }M_j\text{: trace}\ge\text{leading term},
    \label{obc:bulk}
\end{align}
and dividing by $\lambda_j^{2M_j}$ proves (iii). The bulk is now
controlled at every scale; the rest of this section concerns the
boundary only.

\subsection{One doubling step for the opposite-field chain}
\label{obc:steps}

Normalizing by $\lambda_j^m$ removes the extensive part of the
partition function, so the recurrence tracks only the overlap and the
remainder. At scale $j$, write $v_j=v_{\beta_j}$ and, for even $m\ge2$,
define the normalized opposite-field partition function, its remainder,
and the boundary remainder ratio by
\begin{align}
    c_j    &=|\langle\psi_j,v_j\rangle|^2,\notag\\
    B_m(j) &=\frac{Y_m(\beta_j)}{\lambda_j^m}
    =c_j+D_m(j)>0,\notag\\
    D_m(j) &=\folded{R_j^{m/2}}=\|R_j^{m/2}v_j\|^2,\qquad
    e_m(j)=\frac{D_m(j)}{B_m(j)}.
    \label{obc:normalized}
\end{align}
Folding only $R_j^{k/2}$, $k\le m$ even, onto each boundary vector
leaves $R_j^{m-k}$ in the middle, and the remainder decays
geometrically with the extra length (\decayglyph[mpored]):
\begin{equation}
    D_m(j)
    =\tnchain[capl={\bar y},capr=y]{R_j^{m-k}}
    \le\|R_j\|^{m-k}D_k(j),
    \qquad y=R_j^{k/2}v_j,\quad D_m(j)\ge0.
    \label{obc:decay}
\end{equation}
In particular $c_j\le B_m(j)\le B_k(j)$; Steps~1 and~2 use this
monotonicity in the length.

Steps~1 and~2 also need the normalized contraction at the doubled
temperature in terms of quantities at $\beta_j$. Two rows of
Eq.~\eqref{obc:capping} stacked in imaginary time do not fuse site by
site, but the full operators multiply:
\begin{equation}
    \obcstack=\obcstackdouble.
    \label{obc:thermal-diagram}
\end{equation}
Its trace, normalized by $Y_m(\beta_j)^2$, is the thermal purity
$\Pt^{\opp}(m,\beta_j)=\Tr\rho_m^{\opp}(\beta_j)^2\in(0,1]$ of the
Gibbs state $\rho_m^{\opp}(\beta)=e^{-\beta(H_m^{\opp}+\kappa mI)}/Y_m(\beta)$.
Substituting $Y_m(\beta_j)=\lambda_j^mB_m(j)$ and
$\lambda_{j+1}=r_j\lambda_j^2$ relates the normalized contractions at
the two temperatures:
\begin{equation}
    B_m(j+1)
    =\Pt^{\opp}(m,\beta_j)\,r_j^{-m}B_m(j)^2.
    \label{obc:bridge}
\end{equation}
The thermal purity is the only factor not yet tracked; it is entry (T)
of the recurrence.

Given the boundary remainder ratio at length $M_j/2$ and the thermal purity at
length $M_j$, both at $\beta_j$, one step to scale $j+1$ moves through
the (length, inverse temperature) plane in a fixed order; the key lists
the three bounds that Steps~1--3 below derive:
\begin{proofdiagram}
    \obcstairs
\end{proofdiagram}
The proved pair sits at $k_{j+1}=M_j$ and $M_{j+1}=2M_j$, so the step
repeats.

\subsubsection{Step 1: Boundary remainder at the doubled inverse temperature}
\label{obc:move-one}

Step~1 must bound the boundary remainder ratio at the doubled
temperature, which Step~2 then uses. The old overlap says nothing about
the new one $c'$. But $c'$ does not depend on the length, so it cancels
from the difference of two lengths at the new temperature, and
Eq.~\eqref{obc:bridge} returns their ratio to scale $j$. Fix even
lengths $m>k\ge2$, assume $e_k(j)<1$, and write primes for scale $j+1$.
Choose $q<1$ with $\|R_{j+1}\|^{m-k}\le q$, so that
Eq.~\eqref{obc:decay} gives $D_m'\le qD_k'$. Subtracting the two
lengths cancels $c'$:
\begin{align}
    B_k'-B_m'
     &=\Bigl[\cancelled{$\tnchain[capl={\bar v'},capr={v'}]{\Pi'}$}
           +\tnchain[capl={\bar v'},capr={v'}]{(R')^k}\Bigr]
    -\Bigl[\cancelled{$\tnchain[capl={\bar v'},capr={v'}]{\Pi'}$}
         +\tnchain[capl={\bar v'},capr={v'}]{(R')^m}\Bigr]
    =D_k'-D_m'.
    \label{obc:cancellation-diagram}
\end{align}
Combining $D_k'=D_m'+B_k'-B_m'$ with $D_m'\le qD_k'$ gives
\begin{equation}
    (1-q)D_m'\le q(B_k'-B_m'),\qquad
    e_m'\le\frac{q}{1-q}\left(\frac{B_k'}{B_m'}-1\right).
    \label{obc:new-defect}
\end{equation}

Equation~\eqref{obc:bridge} at both lengths expresses the
new-temperature ratio in terms of quantities at scale $j$:
\begin{align}
    \frac{B_k'}{B_m'}
     &=\underbrace{\frac{\Pt^{\opp}(k,\beta_j)}{\Pt^{\opp}(m,\beta_j)}}_{\le1/\Pt^{\opp}(m,\beta_j)}
    \underbrace{r_j^{m-k}}_{\le1}
    \underbrace{\left[\frac{B_k(j)}{B_m(j)}\right]^2}_{\le[1-e_k(j)]^{-2}}\notag\\
     &\le\frac{1}{\Pt^{\opp}(m,\beta_j)[1-e_k(j)]^2},
    \label{obc:ratio}
\end{align}
the last since $B_m(j)\ge c_j=[1-e_k(j)]B_k(j)>0$. Once $e_m'<1$
follows, so does $c'=(1-e_m')B_m'>0$.

Equation~\eqref{obc:new-defect} is a product of two small factors, the
decay $q/(1-q)$ over the length difference and the excess of
$B_k'/B_m'$ above one, which is why the boundary remainder ratio squares.
For the induction, both factors must be $O(x)$ once every tracked bound
is at most $x$. Take
$m=M_j$ and $k=M_j/2$ with $\varepsilon_j\le x\le.01$. Property (ii) at
scale $j+1$ gives $\|R_{j+1}\|^{M_j}\le d(\varepsilon_{j+1})\le d(6x^2)$,
so across the half-chain $m-k=M_j/2$ we may take
\begin{equation}
    q=\sqrt{d(\varepsilon_{j+1})}
    \le\sqrt{\frac{6x^2}{2-18x^2}}\le2x.
    \label{obc:half-chain}
\end{equation}
Separately, when $e_k(j)\le x$ and
$\Pt^{\opp}(m,\beta_j)\ge1-x$, Eq.~\eqref{obc:ratio} bounds the excess above one by
\begin{equation}
    \frac{B_k'}{B_m'}-1\le(1-x)^{-3}-1=O(x).
    \label{obc:ratio-excess}
\end{equation}
Multiplying Eqs.~\eqref{obc:half-chain} and~\eqref{obc:ratio-excess}
in Eq.~\eqref{obc:new-defect} bounds the new boundary remainder ratio $e_m'$ by
$O(x^2)$.

\subsubsection{Step 2: Thermal purity at the doubled length}
\label{obc:move-two}

Step~2 must bound the thermal purity at length $2M$ and the old
temperature $\beta_j$, which Step~3 squares. Doubling an open chain is
not a squaring, since the boundary vectors appear once, so no purity
inequality applies directly. Instead, for even $M$, dividing
Eq.~\eqref{obc:bridge} at length $2M$ by that at length $M$ gives three
factors, each already controlled:
\begin{align}
    \frac{\Pt^{\opp}(2M,\beta_j)}{\Pt^{\opp}(M,\beta_j)}
     &=\underbrace{r_j^M}_{\text{(iii)}}\;
    \underbrace{\frac{B_{2M}(j+1)}{B_M(j+1)}}_{\ge1-e_M(j+1)\ \text{(Step 1)}}\;
    \underbrace{\left[\frac{B_M(j)}{B_{2M}(j)}\right]^2}_{\ge1\ \eqref{obc:decay}}\notag\\
     &\ge r_j^M[1-e_M(j+1)].
    \label{obc:length}
\end{align}
The doubling thus costs only the bulk factor of (iii) and the Step~1
ratio.

\subsubsection{Step 3: Thermal squaring at the doubled length}
\label{obc:move-three}

Step~3 returns to the dyadic grid at $(2M_j,\beta_{j+1})$, where the
next scale starts. The open-chain Gibbs state is positive with unit
trace, so the periodic squaring inequality applies to it unchanged; the
boundary plays no role. Squaring the Gibbs state $\rho_{2M}^{\opp}(\beta_j)$,
Eq.~\eqref{obc:thermal-diagram}, and applying the squaring
inequality~\eqref{obc:square} doubles the inverse temperature at length
$2M$:
\begin{equation}
    1-\Pt^{\opp}(2M,\beta_{j+1})
    \le f(1-\Pt^{\opp}(2M,\beta_j)).
    \label{obc:thermal-square}
\end{equation}
With Steps~1 and~2 this closes one scale of the recurrence.

\subsection{Initial rational steps and the small-remainder induction}
\label{obc:initialization}

The three steps of Sec.~\ref{obc:steps} now become a recurrence across
$j\ge2$ on two numbers per scale, iterated from the seeds, with the
comparison lengths $k_j$ of Eq.~\eqref{obc:k}.
Let $E_j$ and $\eta_j$ denote upper bounds on $e_{k_j}(j)$ and
$1-\Pt^{\opp}(M_j,\beta_j)$. Step~1 needs the decay over the length
difference, while (ii) bounds it over the ring length; a fractional
power converts one to the other. One scale $j\to j+1$ updates them by%
~\leanmark{https://sirui-lu.com/CertifiedGaps/docs/CertifiedGaps/Open/Criterion.html\#CertifiedGaps.OpenMomentWitness.open_step}
\begin{align}
    F(q,e,t)   &=\frac{q}{1-q}
    \left[\frac{1}{(1-e)^2t}-1\right],
               &                                       &\text{Step 1, Eqs.~\eqref{obc:new-defect}--\eqref{obc:ratio}},\notag\\
    q_j        &=d(\varepsilon_{j+1})^{(M_j-k_j)/M_j},
               &                                       &\text{(ii) at scale $j+1$, where $n_{j+1}=M_j$},\notag\\
    E_{j+1}    &\ge F(q_j,E_j,1-\eta_j),
               &                                       &\text{Step 1},\notag\\
    A_j        &=(1-\eta_j)(1-\varepsilon_j)
    (1-\varepsilon_{j+1})(1-E_{j+1}),
               &                                       &\text{Step 2, Eq.~\eqref{obc:length}: }\Pt^{\opp}(2M_j,\beta_j)\ge A_j,\notag\\
    \eta_{j+1} &\ge f(1-A_j),
               &                                       &\text{Step 3, Eq.~\eqref{obc:thermal-square}},
    \label{obc:updates}
\end{align}
valid whenever all four bounds lie in $[0,1)$; $F$ increases in $q$
and $e$ and decreases in $t$, so any upper bound on $q_j$ may be
substituted.

The seed bounds $E_2=.14$ and $\eta_2=.21$ are too large for the
analytical induction below, which needs $x\le.01$, so we first run
Eq.~\eqref{obc:updates} for three explicit rational steps:
\begin{center}
    \begin{minipage}[c]{.42\textwidth}
        \centering
        \begin{tabular}{@{}rrrcr@{}}
            \toprule
            $j$ & $E_j$ & $\eta_j$ & $(M_j-k_j)/M_j$ & $q_j\le$ \\
            \midrule
            2   & .14   & .21      & $49/120$        & .142     \\
            3   & .12   & .155     & $1/2$           & .0282    \\
            4   & .016  & .026     & $1/2$           & .0028    \\
            5   & .001  & .001     & ---             & ---      \\
            \bottomrule
        \end{tabular}
    \end{minipage}\hfill
    \begin{minipage}[c]{.55\textwidth}
        \centering
        \obcladder
    \end{minipage}
\end{center}
Following Ref.~\cite{OpenAI2026Boundaryfield}, each entry is rounded
upward to a simple rational. The fractional powers then reduce to
integer power comparisons that a proof assistant checks directly,
$d(\varepsilon_3)^{49}<.142^{120}$ for $q_2$. Each row follows from
the previous one; at $j=2$,
\begin{align*}
    E_3    & :\ \frac{q_2}{1-q_2}\simeq.166,\quad
    \frac{1}{(1-E_2)^2(1-\eta_2)}-1\simeq.712,\quad
    .166\cdot.712\simeq.118\le.12,                                      \\
    \eta_3 & :\ A_2=(1-\eta_2)(1-\varepsilon_2)(1-\varepsilon_3)(1-E_3)
    \simeq.79\cdot.948\cdot.984\cdot.88=.648,\quad
    f(.352)\simeq.147\le.155.
\end{align*}
The first step gains little because $q_2$ is still large; once $q_j$
is small, the bounds square, as the plot above shows. With
$\varepsilon_5=6\varepsilon_4^2<.000015$, the $j=5$ row brings all
three bounds below $.01$, the base case of the small-remainder
induction:
\begin{equation}
    \varepsilon_j,\ e_{k_j}(j),\
    1-\Pt^{\opp}(M_j,\beta_j)\le x_j,\quad j\ge5;\qquad
    x_5=.01,\quad x_{j+1}=10x_j^2,\quad
    x_j=10^{-1-2^{j-5}}.
    \label{obc:induction}
\end{equation}

The induction step of Eq.~\eqref{obc:induction} must return all three
bounds from $x$ to $10x^2$ in one scale. Each step is quadratic in $x$,
but Step~2 consumes the Step~1 output and the ring bound $\varepsilon_j$ enters both
$q$ and $A$, so the constants are tracked jointly. Assume that the
three bounds hold at scale $j\ge5$ and write $x=x_j\le.01$. One step
from scale $j$ to scale $j+1$ ($M_{j+1}=2M_j$) updates the boundary and
thermal bounds in turn:
\begin{proofdiagram}
    \begin{prooflanes}\def\lanedy{1.05}
        \lanebound{b}{e0}{0}{0}{$e_{M_j/2}(j)\le x$}
        \lanebound{t}{t0}{1}{0}{$\Pt^{\opp}(M_j,\beta_j)\ge1-x$}
        \lanebound{b}{e1}{0}{1}{$e_{M_j}(j+1)\le10x^2$}
        \lanebound{to}{t1}{1}{1.5}{$\Pt^{\opp}(2M_j,\beta_j)\ge1-2x-16x^2$}
        \lanebound{t}{t2}{1}{2.5}{$\Pt^{\opp}(2M_j,\beta_{j+1})\ge1-10x^2$}
        \lanemove{e0}{e1}{1}{\eqref{obc:new-defect}}
        \laneuse{t0}{1}
        \lanemove[mpored]{t0}{t1}{2}{\eqref{obc:length}}
        \laneuse{e1}{2}
        \lanemove[mpoblue]{t1}{t2}{3}{square}
    \end{prooflanes}
\end{proofdiagram}
For Step~1, Eqs.~\eqref{obc:half-chain} and~\eqref{obc:ratio-excess}
give $q_j\le2x$ and
\begin{equation}
    F(2x,x,1-x)
    \le\frac{6x^2}{(1-2x)(1-x)^3}<7x^2,
    \label{obc:boundary-step}
\end{equation}
so $E_{j+1}=10x^2$ bounds $e_{M_j}(j+1)$. For Step~2, the four factors
of $A_j$ are at least $1-x$, $1-x$, $1-6x^2$, and $1-10x^2$. Hence
$A_j\ge(1-x)^2(1-16x^2)\ge1-2x-16x^2$, with $16=6+10$ and the positive
cross terms dropped. For Step~3, since $0<x\le.01$, thermal squaring
yields
\begin{equation}
    f(2x+16x^2)<3x^2,
    \label{obc:thermal-step}
\end{equation}
so $\eta_{j+1}=10x^2$ bounds $1-\Pt^{\opp}(2M_j,\beta_{j+1})$, and
$\varepsilon_{j+1}\le6x^2\le10x^2$ closes the induction at scale $j+1$;
the common prefactor $10$ covers the three constants $7$, $3$, and $6$.

\subsection{Equal endpoint fields, odd lengths, and the uniform gap}
\label{obc:comparison}

Equation~\eqref{obc:induction} controls even opposite-field chains at
every $j\ge5$; the theorem needs the thermal purity of odd equal-field
chains, which Eq.~\eqref{obc:target} turns into the gap. Three things
change: the right boundary vector carries the twist $U_P$, the odd
length makes the remainder neither a square nor necessarily positive,
and only even lengths are certified. Define the normalized equal-field
partition function $Q_j(n)=Y_n^{\same}(\beta_j)/\lambda_j^n$. The twist
is harmless for the leading term, because $U_P$ commutes with the
transfer operator and fixes $\psi_j$ by (i), so the overlap $c_j$
survives:
\begin{equation}
    \tnchain[last=sym]{\Pi_j,U_P}=\tnchain{\Pi_j}.
    \label{obc:twist}
\end{equation}
For the remainder, pull half of an even length $M\le n$ onto each
boundary vector, commuting $U_P$ through $R_j^{M/2}$ ($[U_P,R_j]=0$ because
$[U_P,X_\beta]=0$ and $U_P\Pi_j=\Pi_j$); the middle operator then has norm
at most one ($\|R_j\|\le1$, $U_P$ orthogonal):
\begin{align}
    Q_j(n)-c_j
     &=\tnchain[capl={\bar v_j},capr={v_j}]
       {R_j^{M/2},R_j^{n-M}U_P,R_j^{M/2}}\notag\\
     &=\tnchain[capl={\bar y},capr={y}]{R_j^{n-M}U_P},\qquad
    y=R_j^{M/2}v_j.
    \label{obc:signed}
\end{align}
Cauchy--Schwarz bounds the open contraction by the squared norm of its
caps, the even-length opposite-field remainder:
\begin{equation}
    \schwarzfold=D_M(j).
    \label{obc:absolute}
\end{equation}
The odd equal-field chain is thus controlled by the even opposite-field
remainder at any even $M\le n$, and the recurrence certifies $M=M_j$.

Fix a scale $j\ge5$ and an odd length $n$ with $M_j\le n\le2M_j$. The
equal-field purity is a ratio of $Q$ at the two temperatures, as in
Eq.~\eqref{obc:bridge}, so $Q$ needs an upper bound at $\beta_j$ and a
lower bound at $\beta_{j+1}$. Both follow from Eq.~\eqref{obc:absolute}
at the certified even length $M=M_j=k_{j+1}$, where the recurrence
bounds the remainder at both temperatures:
\begin{align}
    Q_j(n)     &\le c_j+D_M(j)=B_M(j),\notag\\
    Q_{j+1}(n) &\ge c_{j+1}-D_M(j+1)\notag\\
               &=[1-2e_M(j+1)]B_M(j+1)\notag\\
               &\ge(1-2x_{j+1})B_M(j+1)>0.
    \label{obc:comparison-bounds}
\end{align}
Equation~\eqref{obc:bridge} at length $M$ and
Eq.~\eqref{obc:comparison-bounds} then bound the equal-field thermal
purity factor by factor; the bulk factor is the only term that still
depends on $n$:
\begin{align}
    \Pt^{\same}(n,\beta_j)
     &=r_j^n\frac{Q_{j+1}(n)}{Q_j(n)^2}\notag\\
     &\ge\underbrace{(1-2x_{j+1})}_{\eqref{obc:comparison-bounds}}\;
    \underbrace{\Pt^{\opp}(M,\beta_j)}_{\ge1-x_j}\;
    \underbrace{r_j^{n-M}}_{\ge r_j^{M},\ n-M\le M}\notag\\
     &\ge(1-2x_{j+1})(1-x_j)
    \underbrace{(1-\varepsilon_j)(1-\varepsilon_{j+1})}_{\text{(iii)}}\notag\\
     &\ge1-2x_j-26x_j^2\notag\\
     &\ge1-4x_j,
    \label{obc:final-purity}
\end{align}
where $26=20+6$ collects $2x_{j+1}=20x_j^2$ and
$\varepsilon_{j+1}\le6x_j^2$, and $26x_j^2\le2x_j$ for $x_j\le1/13$.
Every odd $n$ in $[M_j,2M_j]$ thus inherits the even-length bound at
scale $j$ up to a factor $4$, which the proof below converts into the gap.

\begin{proof}[Proof of Theorem~\ref{obc:theorem}]
    Each odd $n$ gets its own scale $j$, so the gap bound must be uniform
    in $j$. Both $-\log(1-\Pt^{\same}(n,\beta_j))$ and the inverse temperature
    $\beta_j$ double from one scale to the next, so their ratio in
    Eq.~\eqref{obc:target} is scale independent.
    The intervals $[M_j,2M_j]$, $j\ge5$, cover every odd $n\ge M_5+1=1921$,
    each at its own inverse temperature:
    \begin{proofdiagram}
        \oddcoverage
    \end{proofdiagram}
    Given odd $n$, choose $j\ge5$ with $M_j\le n\le2M_j$.
    Since $x_j\le x_5=.01$, Eq.~\eqref{obc:final-purity} gives
    $w_0\ge\Pt^{\same}(n,\beta_j)\ge1-4x_j>1/2$, so $H_n^{\same}$ has a unique ground state. At this
    length $n$ and finite inverse temperature $\beta_j$,
    Eq.~\eqref{obc:target} gives
    \begin{align}
        e^{-\beta_j\Delta_n}
         &\le\frac{4x_j}{1-4x_j}
         &                               &x_j=10^{-1}\,10^{-2^{j-5}},\quad 1-4x_j\ge.96,\notag\\
         &\le\frac{5}{12}\,10^{-2^{j-5}}
        <10^{-2^{j-5}},
         &                               &\tfrac{.4}{.96}=\tfrac{5}{12}<1.
        \label{obc:gap-final}
    \end{align}
    Taking logarithms and dividing by $\beta_j=\beta_5\,2^{j-5}$ with
    $\beta_5=49\cdot2^3=392$ proves $\Delta_n>\tfrac{\log10}{392}$ for every $j$.
\end{proof}

In Lean~4, we formalize both the general open-chain theorem from a
rational certificate%
~\leanmark{https://sirui-lu.com/CertifiedGaps/docs/CertifiedGaps/Open/Certificate.html\#CertifiedGaps.OpenCertificate.gapGT_of_openCertificate}
and the derivation of the four seeds~\eqref{obc:seeds} from the short
boundary trace table%
~\leanmark{https://sirui-lu.com/CertifiedGaps/docs/CertifiedGaps/Heisenberg/BoundaryInitTable.html\#CertifiedGaps.SpinOneHeisenberg.BoundaryInit.boundaryDefect_le_of_table},
and verify the open-chain integer matrix-product trial contractions on
$m=34,78,142,240$ sites and the nine shortest boundary trace totals
directly in the Lean kernel, leaving only the remaining six short-chain
trace totals as hypotheses%
~\leanmark{https://sirui-lu.com/CertifiedGaps/docs/SeedKernel/Heisenberg/Integers.html\#CertifiedGaps.SpinOneHeisenberg.spin_one_boundary_field_gap_of_hornerTotalsLarge}.
Rerunning the same three rational steps of Eq.~\eqref{obc:updates} from
the four seeds~\eqref{obc:seeds} with four-digit rounding and the ring
constant $5.583$ in place of $6$, and entering the quadratic tail at the
scale-$5$ defect $2.241\times10^{-4}$ instead of $10^{-2}$, sharpens the
uniform bound to $\Delta_n>\log(10^6/2241)/392\simeq0.01556$ for every
odd $n\ge1921$%
~\leanmark{https://sirui-lu.com/CertifiedGaps/docs/CertifiedGaps/Heisenberg/OpenChain.html\#CertifiedGaps.SpinOneHeisenberg.gapGT_of_boundaryFieldInputs_sharp}.
